
Nguồn: \href{http://meddom.org/vai-ky-niem-voi-giao-su-phan-dinh-dieu-gudi/}{Hiệu Minh}

Viện KHVN có hai anh tên vần “iêu”: Viện sỹ Nguyễn Văn Hiệu và Giáo sư Phan Đình Diệu. Một anh là viện trưởng, anh kia là viện phó, cùng giỏi và nổi tiếng khắp thế giới.

Cả hai đã về hưu. Một người huân huy chương đầy ngực, người kia chẳng có cái nào. Nhưng không vì thế mà kém cạnh nhau.


Anh Diệu có con trai là Phan Dương Hiệu. Dân hàn lâm đồn thổi, anh Diệu ghét anh Hiệu, đặt tên con để chửi cho sướng. Thật ra, Dương Hiệu nói ngược là Diệu Hương, tên của anh Diệu và chị vợ là Hương ghép lại, chả liên quan gì đến bác Hiệu.

Ai từng làm việc với anh Hiệu thì phục cũng có mà chê cũng lắm, thậm chí còn gọi đùa là Hiêu, vì hứa mà không làm.

Quân của anh Diệu thì tâm phục, khẩu phục tới 90\%, 10\% còn lại là do tỵ hiềm, ghen ghét, hoặc do lề phải trái. Chưa ai dám gọi anh là Điêu mà dân biết tiếng Pháp gọi anh là Monsieur Dieu – anh Trời.

\begin{center}
\textbf{Nhiệt huyết với khoa học}
\end{center}

Nói về đóng góp của Giáo sư Phan Đình Diệu cho khoa học nước nhà thì báo chí viết chán rồi. Tốt nghiệp tiến sỹ toán lý ở tuổi 30 vào cuối năm 1960 tại Liên Xô, thời đó là ghê răng.

Anh có công lớn trong việc tạo ra ngành Tin học Việt Nam, định hướng phát triển trong nhiều năm. Mình từng nghe các seminar về automat hữu hạn, ngôn ngữ hình thức, lý thuyết mật mã, triết học và toán, hội trường đông nghịt.

Tổng Cua từng viết những chương trình dựa vào thuật toán bảo mật của Giáo sư bằng Pascal vài ngàn lệnh cùng với Vũ Đình Hòa và mấy tay chuyên toán cự phách khác. Anh nói về khóa bảo mật cực phức tạp còn dễ hơn các bạn viết còm. Anh bảo, cái đẹp Toán học nằm ở sự đơn giản và logic. Mình viết blog được là học bài của anh Diệu.

Là người có tâm với khoa học và rất yêu toán và tin, anh bỏ không biết bao công sức để xây dựng ngành này cho nước nhà cùng với những giáo sư đầu ngành như Tạ Quang Bửu, Hoàng  Tụy, Lê Văn Thiêm, Nguyễn Văn Đạo…

Đi sang Pháp, thay vì dành tiền mua xe Peugeot về bán lấy lãi như nhiều người, anh mua rất nhiều linh kiện PC cho Viện. Đến nỗi đám quân thương quá, bảo anh đưa tiền, thay vì mua chip, mấy cha láu cá mua một cái xe gửi cho chị Huơng. Làm khoa học cũng phải có “động cơ”, các lão nói vậy.

Nhớ có lần mưa lụt hầm chứa máy tính ODRA, giáo sư đến thấy thảm cảnh đó mà rơi nước mắt vì biết bao tiền của mới mua được cái máy duy nhất thời đó ở miền Bắc.

Giáo sư nổi tiếng vì những quan điểm chính trị đổi mới vì anh nghiên cứu triết học và toán rất kỹ, hiểu rõ qui luật phát triển. Rất nhiều báo cáo của Đại hội Đảng liên quan đến KHKT thường do anh chấp bút. Tuy vậy, nhiều lời cố vấn quan trọng của người trí thức cũng chỉ dừng trên giá sách.

\begin{center}
\textbf{Những dự đoán tương lai}
\end{center}

Tổng Cua nhớ một số dự đoán tầm quốc gia và cả quốc tế của anh Diệu mà sau đó đã đúng.

Thấy Viện sỹ Nguyễn Văn Hiệu bổ nhiệm cán bộ lung tung, rồi biến viện KH VN thành những công ty nhỏ, Giáo sư đoán, cung cách này sẽ đưa viện khoa học đến sự tàn lụi. Quả thật, “nhân bảo như thần bảo”.

Đọc đề án 112 (tin học hóa quốc gia) năm 2001 có nhiều điểm bất bình thường, anh dự báo, dự án sẽ thất bại và gửi thư khẩn cho Thủ tướng Phan Văn Khải. Quả thật sau 5 năm, ban quản lý ra tòa và hàng nghìn tỷ mất mát.

Có lần đi công tác bên Bulgaria (1986), anh Diệu mang theo cháu Phan Dương Hiệu và bà xã Văn Thị Xuân Hương, em gái của giáo sư Văn Như Cương, toàn vần ương.

Anh kể chuyện năm 1954, tốt nghiệp trung học ở Hà Tĩnh rồi đi bộ ra Hà Nội để học trường ĐH Khoa học. Anh nói là gien của anh là làm toán, không phải để làm quan. Nhưng cuối cùng người ta cứ nhét chức cho anh, khổ thế.

Cuối những năm 1970, cấp trên thấy Giáo sư trẻ và rất giỏi, đề nghị đưa anh vào Đảng để làm cán bộ nguồn, dù khi đó đã làm Viện phó Viện KHVN.

Chị Hương, đảng ủy viên, có trách nhiệm phát triển đảng cho chồng. Tỷ tê về tương lai, rồi lợi ích nếu là người của Đảng thì như thế nào, và khuyên, anh chỉ cần viết đơn là họ kết nạp ngay.

Anh nói, để nghĩ 10 ngày rồi trả lời. Đúng hẹn, anh bảo, Đảng muốn thì cứ kết nạp, cần gì phải viết đơn. Thế là chuyện vào Đảng xếp lại. Quan lộ cũng vì thế mà chẳng hanh thông. Ngang ngạnh kiểu ông đồ xứ Nghệ quả có một không hai.

Nhớ lần đi thang máy lên đỉnh Vitusha của thủ đô Sofia đầy tuyết trắng, thay vì mơ mộng với trời mây, bỗng anh than, với cung cách quản lý của các nước XNCH, thế nào khối này cũng sụp đổ cho mà xem. Lúc đó, Tổng Cua nghe rất choáng, và thấy Giáo sư hơi… phản động.

Anh kể, đã sang Hungaria, CHDC Đức, Ba Lan, Tiệp Khắc, Liên Xô, Mỹ và Tây Âu, nhận xét thế là có lý riêng của mình.

Ba năm sau, bức tường Berlin sụp đổ (1989), kéo theo toàn bộ Đông Âu và Liên Xô tan rã.

Anh là đại biểu Quốc hội lúc hơn 40 tuổi, một trong đại biểu trẻ nhất thời đó. Trong một lần được hỏi về một vị quan cao nhất và ảnh hưởng của ông đối với đất nước, anh Diệu có nói đại ý, vị lãnh đạo ấy rất vĩ đại, đưa cuộc kháng chiến chống Mỹ tới thắng lợi, nhưng vĩ đại hơn, nếu ông từ chức ngay bây giờ.

Lúc đó, vị lãnh đạo kia đang ở thế thượng phong sau 1975 chiến thắng vang dội. Nhưng rồi chính sách tập thể HTX, hợp tỉnh, giá lương tiền quá kém cỏi trong thời bình đã đẩy đất nước đến thảm họa. Người dân chỉ mong ông ấy ra đi.

Mang triết lý chiến tranh vào xây dựng thời bình là không thể, mỗi người có một thời. Anh Diệu từng từ quan về làm nghiên cứu khoa học, góp ý phát triển đất nước theo quan điểm của thời đại mới, cho dù những lời của Giáo sư rất khó nghe đối với nhiều chính trị gia.

Thật tiếc, những đầu óc có tầm vóc quốc tế như Giáo sư Diệu mà ít được sử dụng đúng mục đích. Nước mình quá lãng phí chất xám.

\begin{center}
\textbf{Viện trưởng giỏi việc “nước”, đảm việc nhà}
\end{center}

Viện IOIT có mấy cái nhà ngói hình chữ U. Mỗi lần trời mưa, ếch nhái kêu râm ran. Mình thì tối ngủ bàn cơ quan, ăn cơm tập thể chợ  Ngọc Hà. Gia tài là cái xe đạp Wilga mầu đỏ của Ba Lan, đi vài năm đã hỏng. Quần áo rách dần, tích kê mông, đầu gối, áo sờn. Lốp xe vá chằng vá đụp.

Mặt mũi vêu vao như kẻ chết đói. Đang tuổi ăn tuổi ngủ mà có 13kg gạo/tháng, lương 51 đồng. Mỗi lần về quê, mẹ phải cho thêm vài cân gạo cứu đói vào chiếc ruột tượng, gói cho mấy quả trứng và chai mắm tép. Bố thở dài, kỹ sư chó gì mà nghèo rớt mồng tơi. Hay là về đi cày, ăn no hơn. Học cho lắm chữ mà chẳng nên cơm cháo gì.

Thời đó ai cũng thế. Giáo sư Diệu cũng chẳng hơn gì. Phòng Viện trưởng khoảng $8m^2$, đủ kê cái bàn, 4 cái ghế mây, bàn uống nước, sàn lát gạch thô, trần cót ép, không điều hòa, có cái quạt máy điện cơ chạy lừ đừ như ốm đói.

Nhưng đám trẻ cũng vui vì có em Ngân chân dài, xinh đẹp làm thư ký, đi lại phấp phới. Em Thoa má bánh đúc, trông mỡ màng, làm thư viện. Em Liên cá vàng xinh như mộng, đục lỗ bìa máy tính ODRA.

Các viện sỹ trẻ thi nhau đong đưa, gọi viện là Tán tỉnh và Tiêu khiển cũng không ngoa. Anh cũng thông cảm và bảo, đàn ông thấy người đẹp mà không ngắm thì hoặc anh ta nói dối hoặc là bị bệnh.

Anh Diệu ở tầng 5 bên khu lắp ghép Giảng Võ, vì làm quan được lên tầng cao sạch sẽ và mát. Chỉ tội không có nước. Sáng sáng, trước khi đi làm, Giáo sư cởi trần, xách nước từ tầng một lên tầng năm. Thời đó, ban ngày cả nước lo việc nhà, ban đêm cả nhà lo việc nước, Viện trưởng IOIT cũng vậy thôi.

Khi được cấp nhà rộng hơn ở Đồng Xa, mấy phòng liền, cả gia đình chuyển về đó. Thực không ngờ, nơi đó đồng không mông quạnh, chỉ có đồng lúa, ếch nhái kêu, cào cào châu chấu, thiêu thân và bướm đêm bay rợp trời. Điện mất, nước mất, khổ đủ điều. Viện trưởng tiếp tục thức đêm xếp hàng nước và giặt tã.

Là quân của anh ấy mà Tổng Cua chẳng bao giờ biếu xén thủ trưởng. Trong khi mình sang chơi là chị Hương mời vồn vã, chú ở lại ăn cơm, độc thân trông đói hom hem, khổ thân chú. Thủ trưởng hối lộ quân, chuyện ngược đời.

\begin{center}
\textbf{Nghiên cứu bỏ thuốc lá…ra con trai}
\end{center}

Cái nghề nghiên cứu khoa học nước mình hay lắm. Nghiên cứu một đằng, kết quả ra một nẻo. Giáo sư Diệu nghiên cứu cai thuốc lá thì vợ đẻ con trai.

Chả là Giáo sư có hai cô con gái là Quỳnh Dương và Hà Dương đều rất xinh và hiền. Như nhiều nhà, hình như anh mong có con trai.

% \begin{center}
% \includegraphics[scale=0.5]{ntngochai_phandinhdieu}\\
%  Hai ông bà Hương-Diệu. Ảnh: Ngọc Hải
% \end{center}

Anh Diệu nghiện thuốc lá nặng, phòng làm việc lúc nào cũng nghi nghút như cái lò gạch của Nam Cao. Ai đi công tác về cũng biếu thuốc lá, từ 555 đến Pall Mall, Malborro, Tam Đảo, Điện Biên, đủ loại tây ta.

Vợ khuyên bỏ thuốc mãi không được nên ra giá, anh bỏ được, em sẽ đẻ con trai. Giáo sư hứa đại. Thật bất ngờ, cuối năm ấy, cháu trai Phan Dương Hiệu ra đời. Thế là đành phải thực hiện lời hứa. 

Mấy ngày đầu Giáo sư không chịu nổi, cứ cầm điếu thuốc lên, hút vài hơi rồi dụi tắt. Dần dần anh mang các bao thuốc bóc sẵn để khắp nơi trong nhà, trên giá sách, chỗ bàn ăn, văn phòng, nơi hay sờ tới nhất. Nhìn thấy thuốc nhưng không cầm.

Sau ba tuần anh cai hẳn. Vợ đùa dù hơi lo “Anh Diệu bỏ được thuốc nghĩa là anh ấy có thể bỏ vợ được”.

\begin{center}
\textbf{Sửa gáy Giáo sư được trả công bằng…việc làm}
\end{center}
Năm 1977, anh Diệu, lúc đó khoảng 40 tuổi, Viện trưởng Viện Tính toán và Điều khiển đi thăm các nước XHCN và tìm nhân tài cho viện.

Anh tới Ba Lan vào mùa xuân , tuyết vẫn còn dầy. Gần 200 lưu học sinh VN đón giáo sư trẻ ở ký túc xá trên đường $\dot{Z}wirki i Wigury$, đại lộ huyết mạch từ sân bay $Warszawa-Okecie$ về trung tâm Thủ đô Ba Lan.

Bác đoàn trưởng tên là Loan, giám đốc Nha khí tượng Thủy văn đang làm NCS, gõ cửa phòng và bảo “Anh Cua sang nhờ tý”. Mình sang thì gặp Giáo sư Diệu rất trẻ và dễ mến. Anh cười và bảo “Nghe nói cậu biết cắt tóc, chiều nay tớ có cuộc nói chuyện, đi nhiều nước mà không có thời gian. Cậu giúp tớ nhé”.

Giọng Hà Tĩnh hơi nặng, anh xưng cậu cậu tớ tớ, nghe rất thân mật và không hề quan cách. Mới gặp đã thấy rất cảm tình. Trán rộng, mặt sáng, trông rõ là… bác học.

Thời đó sinh viên Ba Lan nổi tiếng tóc dài quần loe, hippi, sexy, thu thập ảnh các cô trần truồng, treo cả lịch 365 thế làm tình trong phòng, chỉ tội không có giáo cụ trực quan để thực hành. Cắt tóc chỉ là xén bớt phần dài dưới gáy.

Trong khi anh Diệu muốn có mái, rẽ ngôi. Loay hoay một hồi, mình cũng sửa xong cái gáy Giáo sư. Anh hỏi Cua “Cậu học nghề chi?”. Mình bảo là Toán Tin. Thế à, cậu về chỗ mình nhé.

Mình bảo, em học dốt lắm, không có bằng đỏ đen gì đâu, lại còn đúp nữa đó. Anh khuyên, tớ không cần bằng đỏ, mà lấy người về nghiên cứu khoa học, và làm được việc. Rồi anh giới thiệu viện, cho số điện thoại, địa chỉ.

Đang du học nên chả nghĩ đến xin việc vì hồi đó cho rằng làm đâu cũng là xây dựng CNXH. Khi về Hà Nội mới hiểu là cái bán kính 3km lấy tâm là bờ Hồ vô cùng quan trọng. Mình lên Nghĩa Đô tìm, không hy vọng anh nhớ mình là ai. Hóa ra Giáo sư nhớ cả tên.

Anh viết thư tay lên Bộ ĐH, thế là mình nghiễm nhiên thành viện sỹ, nằm bàn, cơm tập thể, ước mơ xây XHCN và tiến lên CNCS vào cuối thế kỷ 20.

\begin{center}
\textbf{Khử virus PC được vợ}
\end{center}

Nếu so sánh hai gia đình các anh vần “iệu” ở Viện KH Việt Nam, phải công nhận thuyết “nhân quả” rất đúng. Cha mẹ sống nhân nghĩa để đức cho con. Không biết anh Nguyễn Văn Hiệu thế nào chứ gia đình anh chị Diệu-Hương hạnh phúc, con cái thành đạt.

Con gái thứ là Phan Thị Hà Dương, từng giành được Huy chương Đồng Olympic Toán quốc tế năm 1990, Tiến sỹ Toán khi 26 tuổi và hiện đang công tác trong nước.


% \begin{figure}[!ht]
% \centering
% \includegraphics[scale=0.5]{haduong}
% \caption{Tiến sỹ toán Hà Dương }
% \end{figure}


Cháu Phan Dương Hiệu hồi đi Bulgaria mới 8 tuổi. Lần đi chơi mình bảo, cháu ghi nhật ký thì nên đếm số cột điện treo thang máy từ chân lên đỉnh Vitusha (Sofia) làm kỷ niệm. Có lẽ do gien toán của bố nên cháu thản nhiên, cần gì phải đếm, số cột đã đánh rồi, chỉ cần lên đỉnh và xem số cột cuối là bao nhiêu. Cháu Dương Hiệu hiện là Tiến sỹ Toán đang ở bên Pháp.

Con gái đầu lòng là Quỳnh Dương, đẹp hiền, đảm, ít nói, ngoan, học giỏi và an phận. Hồi ở Đồng Xa, có lần chị Hương nhờ lão Cua tìm thầy trẻ dạy đàn ghi ta cho cháu Quỳnh. Mình hứa lung tung rồi quên tiệt.

Chị Hương quay sang tìm kỹ sư sửa máy tính cho gia đình. Cu Thắng, viện IOIT, khá đẹp trai, có ria mép, thông minh, nhận đến sửa đĩa cứng, rồi cài virus vào, thỉnh thoảng máy lại treo. Mà PC không treo thì Quỳnh Dương cũng gọi anh Thắng đến xem lại cho chắc.

Sau vài tháng, mình hỏi chị Hương là thế nào rồi. Chị bảo chán lắm, tình hình là rất tình hình, lần nào về cũng thấy hai đứa ngồi đối diện nhau, chả thấy nói gì.

Năm ấy Quỳnh Dương sang Pháp du học. Anh chị thở dài, nông nỗi này thì em Quỳnh ế mất. Cả nhà ra sân bay tiễn con, vừa thương vừa nhớ. Cậu Thắng xin đi theo, nhỡ máy laptop của Quỳnh có hỏng gì chăng.

Lên sân bay Nội Bài bịn rịn chia tay. Bỗng anh Diệu thấy con gái cưng, vốn hiền lành, ôm chầm thằng cu chuyên cài virus, nắm tay, nhìn nhau đắm đuối chẳng biết trời đất nào nữa, rồi hôn môi một cái rất dài.

Giáo sư bỏ kính xuống, không thể tin vào mắt mình. Đời anh từng quản lý hàng ngàn cán bộ khoa học, đưa ra những triết lý phát triển cho quốc gia, báo cáo khoa học quốc tế được cả ngàn người thán phục. Rồi suýt bị “dính chưởng” vì chị Dương Thu Hương từng có bài viết đa nguyên của Giáo sư trong vali khi bị bắt ở sân bay. 

Ngay cả cái lý thuyết automat hữu hạn phải gió hay cái máy vi tính đầu tiên của Việt Nam được xuất xưởng, anh cũng không ngạc nhiên bằng nụ hôn giữa thanh thiên bạch nhật của cô con gái  dành cho cu Thắng bên phòng Kỹ thuật số.

Hai vợ chồng choáng váng và chợt hiểu thế nào là virus PC nhiễm sang người. Họ thầm gọi cu Thắng là Thang Dieu – Thằng  Con Trời cho xứng với Monsieur Dieu – Ông Trời Diệu.

Sau cái vụ nhiễm virus ấy, hai đứa có mấy con và hiện đang sống bên Pháp, rất hạnh phúc. Chả hiểu cu Thắng có dạy đàn ghi ta cho Quỳnh không.

Còn Tổng Cua đang làm gì, đố các bạn biết? Lão cố dạy hai thằng cu Hiệu – Minh cách cài virus vào máy tính và học thêm nghề cắt tóc.

HM. 1-2-2012

